%% =============================================================================
%% sm_a.tex — Supplementary Material A: Mathematical Properties of the HBB
%% Body content only (subsections); \section and \label defined in osm_body.tex
%% Primary source: dev/log/03_research-notes-deep-mathematical-properties.qmd
%% =============================================================================

This appendix collects the complete proofs and supporting results for the
mathematical properties of the hurdle beta-binomial model discussed in
\cref{sec:model}. The beta-binomial distribution in the mean--precision
parameterization follows~\citet{Griffiths1973} and~\citet{Prentice1986};
the hurdle construction follows~\citet{Mullahy1986} and~\citet{Cragg1971}.
We maintain the notation established in~\cref{sec:notation}
throughout.

%% ---------------------------------------------------------------------------
%% A.1  Proof of Theorem 2.1 (Monotonicity of h)
%% Source: QMD lines 64--138
%% ---------------------------------------------------------------------------

\subsection{Proof of \texorpdfstring{\cref{thm:monotonicity}}{Theorem~2.1}}
\label{sma:proof-monotonicity}

We give the complete proof that the intensity function
$h(\mu) = \mu / (1 - p_0)$ is strictly increasing in $\mu$ for every
$n \ge 2$, $\kappa > 0$. Recall that
$p_0(n,\mu,\kappa) = \prod_{j=0}^{n-1}[(1-\mu)\kappa + j]/(\kappa + j)$
is the zero probability of the beta-binomial distribution, as defined
in~\eqref{eq:p0}.

\paragraph{Auxiliary quantities.}
Write $b = (1-\mu)\kappa$ and define
\begin{equation}\label{sma:eq-Lambda}
  \Lambda
  = \kappa \sum_{j=0}^{n-1} \frac{1}{(1-\mu)\kappa + j}
  > 0,
  \qquad
  \Lambda_2
  = \kappa^2 \sum_{j=0}^{n-1} \frac{1}{\bigl[(1-\mu)\kappa + j\bigr]^2}
  > 0.
\end{equation}
Note that $\partial p_0/\partial\mu = -p_0\,\Lambda$ (obtained by
differentiating $\ln p_0 = \sum_{j=0}^{n-1}\ln[(1-\mu)\kappa + j]
- \ln[\kappa + j]$) and $\partial\Lambda/\partial\mu = \Lambda_2$
(since $\partial b/\partial\mu = -\kappa$).

\begin{lemma}[Sum-of-reciprocals inequality]\label{sma:lem-reciprocals}
  Let $b > 0$, $n \ge 1$, and set
  $S_1 = \sum_{j=0}^{n-1}(b+j)^{-1}$ and
  $S_2 = \sum_{j=0}^{n-1}(b+j)^{-2}$, so that
  $\Lambda = \kappa\,S_1$ and $\Lambda_2 = \kappa^2\,S_2$.
  \begin{enumerate}[label=\textup{(\roman*)}]
    \item For $n = 1$: $S_1^2 = S_2$, so $\Lambda^2 = \Lambda_2$.
    \item For $n \ge 2$: $S_1^2 > S_2$, so $\Lambda^2 > \Lambda_2$.
  \end{enumerate}
\end{lemma}

\begin{proof}
  Write $a_j = (b + j)^{-1}$.  Expanding the square of the sum,
  \[
    S_1^2
    = \Bigl(\sum_{j=0}^{n-1} a_j\Bigr)^{\!2}
    = \sum_{j=0}^{n-1} a_j^2
      + 2\!\!\sum_{0 \le i < j \le n-1}\!\! a_i\, a_j
    = S_2
      + 2\!\!\sum_{0 \le i < j \le n-1}\!\!
        \frac{1}{(b+i)(b+j)}.
  \]
  For $n = 1$ the cross-product sum is empty, giving $S_1^2 = S_2$.
  For $n \ge 2$ every term $1/[(b+i)(b+j)]$ is strictly positive, so
  $S_1^2 > S_2$.
\end{proof}

\begin{proof}[Proof of~\cref{thm:monotonicity}]
  Define the numerator function
  \begin{equation}\label{sma:eq-Phi}
    \Phi(\mu)
    = \bigl(1 - p_0(\mu)\bigr) - \mu\, p_0(\mu)\,\Lambda(\mu).
  \end{equation}
  From the quotient rule applied to $h = \mu\,(1-p_0)^{-1}$ together with
  $\partial p_0 / \partial\mu = -p_0\, \Lambda$, one obtains
  \begin{equation}\label{sma:eq-dh-numerator}
    \frac{\partial h}{\partial \mu}
    = \frac{\Phi(\mu)}{(1 - p_0)^2}.
  \end{equation}
  Since $(1 - p_0)^2 > 0$ for $\mu \in (0,1)$, the sign of
  $\partial h / \partial\mu$ equals the sign of $\Phi(\mu)$.  The proof
  proceeds in three steps.

  \medskip\noindent\textit{Step~1: Boundary value.}\enspace
  At $\mu = 0$, every factor in~\eqref{eq:p0} equals one, so $p_0(0) = 1$.
  Hence $\Phi(0) = (1 - 1) - 0 \cdot 1 \cdot \Lambda(0) = 0$.

  \medskip\noindent\textit{Step~2: Derivative of\/ $\Phi$.}\enspace
  Differentiating~\eqref{sma:eq-Phi} with the product rule gives
  \[
    \Phi'(\mu)
    = p_0\,\Lambda
      - \bigl[p_0\,\Lambda
        + \mu\,\tfrac{d}{d\mu}(p_0\,\Lambda)\bigr]
    = -\mu\,\frac{d}{d\mu}\bigl(p_0\,\Lambda\bigr).
  \]
  Now $d(p_0\,\Lambda)/d\mu
  = (\partial p_0/\partial\mu)\,\Lambda + p_0\,(\partial\Lambda/\partial\mu)
  = -p_0\,\Lambda^2 + p_0\,\Lambda_2
  = p_0(\Lambda_2 - \Lambda^2)$,
  and therefore
  \begin{equation}\label{sma:eq-Phi-deriv}
    \Phi'(\mu)
    = -\mu\, p_0\bigl(\Lambda_2 - \Lambda^2\bigr)
    = \mu\, p_0\bigl(\Lambda^2 - \Lambda_2\bigr).
  \end{equation}

  \medskip\noindent\textit{Step~3: Application of
  \cref{sma:lem-reciprocals}.}\enspace
  \begin{itemize}[nosep]
    \item \emph{Case $n = 1$.}
      \cref{sma:lem-reciprocals}(i) gives $\Lambda^2 = \Lambda_2$, so
      $\Phi'(\mu) = 0$ identically.  Combined with $\Phi(0) = 0$, we obtain
      $\Phi(\mu) = 0$ for all $\mu \in (0,1)$ and hence
      $\partial h / \partial\mu = 0$.  Equivalently, $p_0 = 1 - \mu$ when
      $n = 1$, so $h(\mu) = \mu / (1 - (1 - \mu)) = 1$ identically.
    \item \emph{Case $n \ge 2$.}
      \cref{sma:lem-reciprocals}(ii) gives $\Lambda^2 > \Lambda_2$.
      Since $\mu > 0$ and $p_0 > 0$ on $(0,1)$,
      equation~\eqref{sma:eq-Phi-deriv} yields $\Phi'(\mu) > 0$ for all
      $\mu \in (0,1)$.
  \end{itemize}
  Since $\Phi(0) = 0$ and $\Phi$ is strictly increasing on $(0,1)$, we
  conclude $\Phi(\mu) > 0$ for all $\mu \in (0,1)$, and therefore
  \[
    \frac{\partial h}{\partial\mu}
    = \frac{\Phi(\mu)}{(1 - p_0)^2} > 0. \qedhere
  \]
\end{proof}

%% ---------------------------------------------------------------------------
%% A.2  Higher Moments of the Truncated Beta-Binomial
%% Source: QMD lines 141--260
%% ---------------------------------------------------------------------------

\subsection{Higher moments of the truncated beta-binomial}
\label{sma:higher-moments}

\paragraph{Notation.}
Throughout this subsection, $Y \sim \BetaBin(n,\mu,\kappa)$ denotes the
standard (untruncated) beta-binomial with $a = \mu\kappa$,
$b = (1-\mu)\kappa$, and
$Y^+ \sim \ZTBB(n,\mu,\kappa)$ denotes the zero-truncated version.  We
write $p_0 = \Prob(Y = 0)$,
$V_{\mathrm{BB}} = \Var(Y) = n\mu(1-\mu)(n+\kappa)/(1+\kappa)$, and
$h = \mu/(1-p_0)$.

%% --- Second factorial moment ---

\begin{lemma}[Second factorial moment]\label{sma:lem-second-factorial}
  For $Y \sim \BetaBin(n,\mu,\kappa)$,
  \begin{equation}\label{sma:eq-second-factorial}
    \boxed{\E\bk{Y(Y-1)} = n(n-1) \cdot \frac{\mu(\mu\kappa + 1)}{\kappa + 1}.}
  \end{equation}
\end{lemma}

\begin{proof}
  The beta-binomial arises as a beta
  mixture~\citep{Griffiths1973}:
  let $p \sim \Beta(a,b)$ with $a = \mu\kappa$,
  $b = (1\!-\!\mu)\kappa$, and $Y \mid p \sim \Bin(n,p)$.
  Since $\E[Y(Y\!-\!1) \mid p] = n(n\!-\!1)p^2$,
  the tower property gives
  \[
    \E\bk{Y(Y-1)} = n(n-1)\,\E[p^2].
  \]
  The second moment of $p \sim \Beta(a,b)$ is
  $\E[p^2] = a(a+1)/[(a+b)(a+b+1)]$, using the rising factorial identity
  $\E[p^k] = a^{(k)}/(a+b)^{(k)}$ for $k = 2$.  Substituting
  $a = \mu\kappa$, $a + b = \kappa$:
  \[
    \E\bk{Y(Y-1)}
    = n(n-1) \cdot \frac{\mu\kappa(\mu\kappa + 1)}{\kappa(\kappa+1)}
    = n(n-1) \cdot \frac{\mu(\mu\kappa + 1)}{\kappa+1}. \qedhere
  \]
\end{proof}

\paragraph{Consistency check.}
From~\eqref{sma:eq-second-factorial}, the second moment is
$\E[Y^2] = \E[Y(Y-1)] + \E[Y] = n(n-1)\mu(\mu\kappa+1)/(\kappa+1) + n\mu$.
We confirm $\Var(Y) = \E[Y^2] - (n\mu)^2$:
\[
  \E[Y^2] - n^2\mu^2
  = \frac{n(n-1)\mu(\mu\kappa+1) + n\mu(\kappa+1) - n^2\mu^2(\kappa+1)}
         {\kappa+1}.
\]
Expanding the numerator factor:
$n\mu\bk{(n-1)(\mu\kappa+1) + (\kappa+1) - n\mu(\kappa+1)}$.  The bracket
simplifies as
\[
  (n-1)\mu\kappa + (n-1) + (\kappa+1) - n\mu\kappa - n\mu
  = (1-\mu)(n + \kappa).
\]
Therefore
$\Var(Y) = n\mu(1-\mu)(n+\kappa)/(\kappa+1) = V_{\mathrm{BB}}$.
\enspace$\square$

%% --- General r-th factorial moment ---

\begin{lemma}[General $r$-th factorial moment]\label{sma:lem-rth-factorial}
  For $Y \sim \BetaBin(n,\mu,\kappa)$,
  \begin{equation}\label{sma:eq-rth-factorial}
    \E\!\left[\prod_{j=0}^{r-1}(Y-j)\right]
    = \frac{n^{(r)} \cdot (\mu\kappa)^{(r)}}{\kappa^{(r)}},
  \end{equation}
  where $c^{(r)} = c(c+1)\cdots(c+r-1) = \Gamma(c+r)/\Gamma(c)$ is the
  rising factorial.
\end{lemma}

\begin{proof}
  Follows from
  $\E[\prod_{j=0}^{r-1}(Y-j) \mid p] = n^{(r)}p^r$ for the binomial, and
  $\E[p^r] = (\mu\kappa)^{(r)}/\kappa^{(r)}$ for the beta.
\end{proof}

%% --- Truncation identity ---

\begin{proposition}[Truncation identity for functions vanishing at zero]
  \label{sma:prop-trunc-identity}
  For any function $g$ with $g(0) = 0$,
  \begin{equation}\label{sma:eq-trunc-identity}
    \E\bk{g(Y) \mid Y > 0} = \frac{\E\bk{g(Y)}}{1 - p_0}.
  \end{equation}
  In particular,
  $\E[Y^k \mid Y > 0] = \E[Y^k]/(1-p_0)$ for all $k \ge 1$.
\end{proposition}

\begin{proof}
  Since $g(Y) \cdot \one(Y = 0) = g(0) \cdot \one(Y = 0) = 0$, we have
  $\E[g(Y)] = \E[g(Y) \cdot \one(Y > 0)]$.  Dividing by
  $\Prob(Y > 0) = 1 - p_0$ gives the result.
\end{proof}

%% --- Conditional second factorial moment ---

\begin{corollary}[Conditional second factorial moment]
  \label{sma:cor-cond-second-factorial}
  \begin{equation}\label{sma:eq-cond-second-factorial}
    \boxed{\E\bk{Y(Y-1) \mid Y > 0}
    = \frac{n(n-1)\mu(\mu\kappa+1)}{(\kappa+1)(1-p_0)}.}
  \end{equation}
\end{corollary}

%% --- Truncated BB variance ---

\begin{proposition}[Truncated BB variance]
  \label{sma:prop-trunc-var}
  For $Y^+ \sim \ZTBB(n,\mu,\kappa)$,
  \begin{equation}\label{sma:eq-trunc-var}
    \boxed{\Var(Y \mid Y > 0)
    = \frac{V_{\mathrm{BB}}(1-p_0) - n^2\mu^2\,p_0}{(1-p_0)^2}.}
  \end{equation}
\end{proposition}

\begin{proof}
  By the truncation identity (\cref{sma:prop-trunc-identity}),
  \[
    \E[Y^2 \mid Y > 0] = \frac{\E[Y^2]}{1 - p_0},
    \qquad
    \E[Y \mid Y > 0] = \frac{n\mu}{1 - p_0}.
  \]
  Then
  \[
    \Var(Y \mid Y > 0)
    = \frac{\E[Y^2]}{1-p_0} - \frac{n^2\mu^2}{(1-p_0)^2}
    = \frac{V_{\mathrm{BB}} + n^2\mu^2}{1-p_0}
      - \frac{n^2\mu^2}{(1-p_0)^2}.
  \]
  Combining the $n^2\mu^2$ terms over the common denominator $(1-p_0)^2$:
  \[
    = \frac{V_{\mathrm{BB}}}{1-p_0}
      + n^2\mu^2 \cdot \frac{(1-p_0) - 1}{(1-p_0)^2}
    = \frac{V_{\mathrm{BB}}}{1-p_0}
      - \frac{n^2\mu^2\,p_0}{(1-p_0)^2}
    = \frac{V_{\mathrm{BB}}(1-p_0) - n^2\mu^2\,p_0}{(1-p_0)^2}.
    \qedhere
  \]
\end{proof}

\begin{remark}\label{sma:rem-competing-effects}
  The formula~\eqref{sma:eq-trunc-var} reveals two competing effects of
  truncation.  The first term $V_{\mathrm{BB}}/(1-p_0)$ inflates the
  variance by rescaling probabilities.  The second term
  $n^2\mu^2 p_0/(1-p_0)^2$ reduces the variance because the truncated
  distribution loses mass at zero, pulling the mean away from zero and
  tightening the effective spread relative to that shifted mean.
\end{remark}

\paragraph{Alternative forms.}

\emph{Form~A (factored):}
\begin{equation}\label{sma:eq-trunc-var-formA}
  \Var(Y \mid Y > 0)
  = \frac{n\mu}{1-p_0}
    \left[\frac{(1-\mu)(n+\kappa)}{1+\kappa}
          - \frac{n\mu\,p_0}{1-p_0}\right].
\end{equation}
This displays the variance as $\E[Y \mid Y > 0]$ times a ``spread factor.''

\emph{Form~B (via factorial moments):}
\begin{equation}\label{sma:eq-trunc-var-formB}
  \Var(Y \mid Y > 0)
  = \frac{\E[Y(Y-1)]}{1-p_0}
    + \frac{n\mu}{1-p_0}
    - \frac{n^2\mu^2}{(1-p_0)^2}.
\end{equation}

\paragraph{Verification of limiting cases.}

\emph{Case~1: $\kappa \to \infty$ (truncated binomial limit).}
As $\kappa \to \infty$, $V_{\mathrm{BB}} \to n\mu(1-\mu)$ and
$p_0 \to (1-\mu)^n$.  The formula becomes the known variance of the
zero-truncated binomial.  For $n = 1$: $Y \mid Y > 0 = 1$
deterministically, so $\Var = 0$.  In the formula: numerator
${}= \mu(1-\mu) \cdot \mu - \mu^2(1-\mu) = 0$.\enspace$\square$

The remaining limiting cases are verified analogously:
as $p_0 \to 0$ the variance reduces to $V_{\mathrm{BB}}$, and
for $n=1$ both sides equal zero since $Y \mid Y > 0 = 1$
deterministically.\enspace$\square$

\paragraph{Numerical cross-check.}
Consider $n = 5$, $\mu = 0.3$, $\kappa = 2$ (so $a = 0.6$, $b = 1.4$).
Direct computation gives
$p_0 = \prod_{j=0}^{4}(1.4+j)/(2+j) = 0.376992$,
$\E[Y] = 1.5$,
$V_{\mathrm{BB}} = 2.45$, and
$\E[Y^2] = 4.70$.
The truncated moments are
$\E[Y \mid Y > 0] = 1.5/0.623008 = 2.40770$,
$\E[Y^2 \mid Y > 0] = 4.70/0.623008 = 7.54409$, and
$\Var(Y \mid Y > 0) = 7.54409 - 2.40770^2 = 1.74707$.
Form~A gives
$(2.45 \times 0.623008 - 2.25 \times 0.376992)/0.623008^2 = 1.74709$,
agreeing to five significant figures.\enspace$\square$

%% ---------------------------------------------------------------------------
%% A.3  Unconditional variance decomposition
%% Source: QMD lines 263--360 (03_research-notes-deep-mathematical-properties)
%% Corrections applied: T1-3 (exact OD formula), Erratum at QMD line 321
%% ---------------------------------------------------------------------------

\subsection{Unconditional variance decomposition}\label{sma:vardecomp}

We derive the unconditional variance of $Y_i$ under the hurdle
beta-binomial model, its variance-to-mean ratio, an overdispersion index
relative to the binomial, and the coefficient of variation. Throughout, we
write $V_{\textup{BB},i} = n_i\mu_i(1-\mu_i)(n_i + \kappa)/(1 + \kappa)$
for the beta-binomial variance~\eqref{eq:bb-moments}, $h_i = \mu_i/(1-p_{0,i})$
for the intensity function~\eqref{eq:uncond-mean}, and
$m_i^{+} = n_i h_i = n_i\mu_i/(1-p_{0,i})$ for the conditional mean given
participation.

%% --- Law of total variance ---

\paragraph{Law of total variance.}

\begin{proposition}[Hurdle model variance decomposition]
\label{sma:prop-var-decomp}
Let $Y_i$ follow the $\HBB$ model
\citep[cf.][for the Poisson hurdle analogue]{Mullahy1986}
with participation probability~$q_i$ and
intensity parameters $(n_i, \mu_i, \kappa)$. Define
$Z_i = \one\{Y_i > 0\}$. Then
\begin{equation}\label{sma:eq-var-decomp}
  \Var(Y_i)
  = q_i \cdot \Var(Y_i \mid Y_i > 0)
    + q_i(1-q_i)\bigl[\E(Y_i \mid Y_i > 0)\bigr]^2.
\end{equation}
\end{proposition}

\begin{proof}
Apply the law of total variance with respect to the Bernoulli variable~$Z_i$:
\[
  \Var(Y_i)
  = \E\bigl[\Var(Y_i \mid Z_i)\bigr]
    + \Var\bigl(\E[Y_i \mid Z_i]\bigr).
\]
\emph{First term.}  When $Z_i = 0$, $Y_i = 0$ almost surely, so
$\Var(Y_i \mid Z_i = 0) = 0$.  When $Z_i = 1$,
$\Var(Y_i \mid Z_i = 1) = \Var(Y_i \mid Y_i > 0)$.  Hence
\[
  \E\bigl[\Var(Y_i \mid Z_i)\bigr]
  = q_i\,\Var(Y_i \mid Y_i > 0).
\]
\emph{Second term.}  The conditional expectation $\E[Y_i \mid Z_i]$ takes
value~$0$ with probability $1-q_i$ and $m_i^{+} = n_i\mu_i/(1-p_{0,i})$
with probability~$q_i$.  Its variance is therefore
$q_i(1-q_i)(m_i^{+})^2$.
\end{proof}

\begin{corollary}[Expanded form]\label{sma:cor-var-expanded}
Substituting the truncated variance from
\cref{sma:prop-trunc-var} into~\eqref{sma:eq-var-decomp}:
\begin{equation}\label{sma:eq-var-expanded}
  \Var(Y_i)
  = \frac{q_i\, V_{\textup{BB},i}}{1-p_{0,i}}
    - \frac{q_i\, n_i^2 \mu_i^2\, p_{0,i}}{(1-p_{0,i})^2}
    + \frac{q_i(1-q_i)\, n_i^2 \mu_i^2}{(1-p_{0,i})^2}.
\end{equation}
Combining the terms involving $n_i^2\mu_i^2/(1-p_{0,i})^2$ yields the
compact form
\begin{equation}\label{sma:eq-var-compact}
  \Var(Y_i)
  = \frac{q_i\, V_{\textup{BB},i}}{1-p_{0,i}}
    + \frac{q_i\, n_i^2\mu_i^2\,(1-q_i-p_{0,i})}{(1-p_{0,i})^2}.
\end{equation}
\end{corollary}

\begin{remark}\label{sma:rem-qi-vs-p0}
The quantities $q_i$ and $p_{0,i}$ are distinct: $q_i$ is the hurdle
participation probability, whereas $p_{0,i}$ is the zero probability of the
beta-binomial component. In general
$q_i \neq 1 - p_{0,i}$; the quantity $(1 - q_i - p_{0,i})$ may be positive,
negative, or zero, so the second term in~\eqref{sma:eq-var-compact} can
either inflate or reduce the unconditional variance.
\end{remark}

\begin{corollary}[In terms of the unconditional mean]
\label{sma:cor-var-uncond-mean}
Recall from~\eqref{eq:uncond-mean} that $\E[Y_i] = q_i\,n_i\,h_i$. Then
\begin{equation}\label{sma:eq-var-uncond-mean}
  \Var(Y_i)
  = q_i\,\Var(Y_i \mid Y_i > 0)
    + \frac{1-q_i}{q_i}\,\bigl(\E[Y_i]\bigr)^2.
\end{equation}
The between-component term is proportional to the squared mean, scaled by
the odds of non-participation $(1-q_i)/q_i$.
\end{corollary}

%% --- Variance-to-mean ratio ---

\paragraph{Variance-to-mean ratio.}

\begin{proposition}[Variance ratio]\label{sma:prop-VR}
Define $\textup{VR}_i = \Var(Y_i)/\E[Y_i]$.  Then
\begin{equation}\label{sma:eq-VR}
  \boxed{
    \textup{VR}_i
    = \frac{(1-\mu_i)(n_i+\kappa)}{1+\kappa}
      - \frac{n_i\,\mu_i\,p_{0,i}}{1-p_{0,i}}
      + (1-q_i)\,n_i\,h_i.
  }
\end{equation}
\end{proposition}

\begin{proof}
Using $\E[Y_i] = q_i\,n_i\,h_i$ and~\cref{sma:cor-var-uncond-mean}:
\[
  \textup{VR}_i
  = \frac{\Var(Y_i \mid Y_i > 0)}{n_i\,h_i}
    + (1-q_i)\,n_i\,h_i.
\]
For the first term, substitute the truncated variance
(\cref{sma:prop-trunc-var}) and simplify using
$h_i = \mu_i/(1-p_{0,i})$:
\[
  \frac{V_{\textup{BB},i}}{(1-p_{0,i})\,n_i\,h_i}
  = \frac{(1-\mu_i)(n_i+\kappa)}{1+\kappa},
  \qquad
  \frac{n_i\,\mu_i^2\,p_{0,i}}{(1-p_{0,i})^2\,h_i}
  = \frac{n_i\,\mu_i\,p_{0,i}}{1-p_{0,i}}. \qedhere
\]
\end{proof}

\begin{remark}\label{sma:rem-VR-interpretation}
The three terms in~\eqref{sma:eq-VR} have clear interpretations:
\begin{enumerate}[label=\textup{(\roman*)}]
  \item $(1-\mu_i)(n_i+\kappa)/(1+\kappa)$: the beta-binomial variance
    inflation factor~\citep{Prentice1986}, which exceeds unity for all $n_i > 1$;
  \item $-n_i\,\mu_i\,p_{0,i}/(1-p_{0,i})$: a correction from
    zero-truncation that reduces the variance ratio, reflecting the
    mass removed at $Y = 0$;
  \item $(1-q_i)\,n_i\,h_i$: the between-component contribution from the
    hurdle, which increases the variance ratio whenever $q_i < 1$.
\end{enumerate}
For Poisson data, $\textup{VR} = 1$ (equidispersion). In the NSECE data,
$\textup{VR}_i > 1$ consistently, reflecting both overdispersion and the
hurdle mechanism.
\end{remark}

%% --- Overdispersion index ---

\paragraph{Overdispersion index relative to the binomial.}

\begin{definition}[Overdispersion index]\label{sma:def-OD}
For a proportion $Y_i/n_i$, define
\begin{equation}\label{sma:eq-OD}
  \textup{OD}_i
  = \frac{\Var(Y_i)}{n_i\,\pi_i(1-\pi_i)},
  \qquad
  \text{where }
  \pi_i = q_i\,h_i = \frac{\E[Y_i]}{n_i}.
\end{equation}
\end{definition}

The denominator $n_i\,\pi_i(1-\pi_i)$ is the variance of a
$\Bin(n_i, \pi_i)$ random variable with the same unconditional mean
proportion, so $\textup{OD}_i$ measures overdispersion relative to the
binomial benchmark.

\paragraph{Limiting cases.}

\begin{itemize}[nosep]
  \item \emph{No hurdle, $p_0 \to 0$.} Then $q_i = 1$, $h_i \to \mu_i$,
    and $\Var(Y_i) \to V_{\textup{BB},i}$, so
    $\textup{OD}_i \to (n_i+\kappa)/(1+\kappa)$, the standard beta-binomial
    overdispersion factor~\citep{Prentice1986}.
  \item \emph{Binomial limit ($\kappa \to \infty$, $q_i = 1$).}
    $V_{\textup{BB},i} \to n_i\mu_i(1-\mu_i)$ and the hurdle is inactive, so
    $\textup{OD}_i \to 1$.
\end{itemize}

%% --- Coefficient of variation ---

\paragraph{Coefficient of variation.}

\begin{proposition}[Conditional CV]\label{sma:prop-CV-cond}
The squared coefficient of variation of $Y_i$ conditional on participation
is
\begin{equation}\label{sma:eq-CV-cond}
  \boxed{
    \textup{CV}^2(Y_i \mid Y_i > 0)
    = \frac{(1-\mu_i)(n_i+\kappa)}{n_i\,\mu_i(1+\kappa)}\,(1-p_{0,i})
      - p_{0,i}.
  }
\end{equation}
\end{proposition}

\begin{proof}
By definition,
$\textup{CV}^2 = \Var(Y_i \mid Y_i > 0) / [\E(Y_i \mid Y_i > 0)]^2$.
Substituting the truncated variance (\cref{sma:prop-trunc-var}) and the
conditional mean $m_i^{+} = n_i\mu_i/(1-p_{0,i})$:
\begin{align*}
  \textup{CV}^2
  &= \frac{V_{\textup{BB},i}(1-p_{0,i}) - n_i^2\mu_i^2\,p_{0,i}}
         {n_i^2\mu_i^2}
  = \frac{V_{\textup{BB},i}}{n_i^2\mu_i^2}\,(1-p_{0,i}) - p_{0,i} \\
  &= \frac{(1-\mu_i)(n_i+\kappa)}{n_i\,\mu_i(1+\kappa)}\,(1-p_{0,i})
    - p_{0,i}. \qedhere
\end{align*}
\end{proof}

\begin{proposition}[Unconditional CV]\label{sma:prop-CV-uncond}
The squared coefficient of variation of $Y_i$ (unconditional) satisfies
\begin{equation}\label{sma:eq-CV-uncond}
  \boxed{
    \textup{CV}^2(Y_i)
    = \frac{\textup{CV}^2(Y_i \mid Y_i > 0) + 1 - q_i}{q_i}.
  }
\end{equation}
\end{proposition}

\begin{proof}
From $\Var(Y_i) = q_i\,\Var(Y_i \mid Y_i > 0) + q_i(1-q_i)(n_i h_i)^2$
(\cref{sma:prop-var-decomp}) and $[\E(Y_i)]^2 = q_i^2(n_i h_i)^2$:
\[
  \textup{CV}^2(Y_i)
  = \frac{\textup{CV}^2(Y_i \mid Y_i > 0)}{q_i}
    + \frac{1-q_i}{q_i}
  = \frac{\textup{CV}^2(Y_i \mid Y_i > 0) + 1 - q_i}{q_i}. \qedhere
\]
\end{proof}

\begin{remark}\label{sma:rem-CV-PPC}
The coefficient of variation is useful for posterior predictive checking
because: \textup{(i)}~it is scale-free, enabling comparison across
providers with different~$n_i$; \textup{(ii)}~the additive $(1-q_i)/q_i$
term in~\eqref{sma:eq-CV-uncond} reveals how the hurdle magnifies the CV
when participation is low; and \textup{(iii)}~observed sample CVs can be
compared against the posterior distribution of theoretical CVs, providing a
diagnostic that simultaneously probes the mean, variance, and
zero-inflation structure of the model.
\end{remark}

%% ---------------------------------------------------------------------------
%% A.4  Analytic properties of the intensity function
%% Source: QMD lines 363--524 (03_research-notes-deep-mathematical-properties)
%% Correction applied: Prop 5.10(b) limit ε_h → 1−C_κ (not 1)
%% Issue 20 applied: Prop 5.13 → Conjecture A.1
%% ---------------------------------------------------------------------------

\subsection{Analytic properties of the intensity function}\label{sma:intensity}

This subsection collects analytic results for the intensity function
$h(\mu) = \mu/(1-p_0)$ that underpin the elasticity formula
(\cref{prop:elasticity}), the boundary behavior discussed in
\cref{rem:boundary}, and the log-scale decomposition of marginal effects.
We continue to write $b = (1-\mu)\kappa$ and use
$\Lambda$, $\Lambda_2$ as defined in~\eqref{sma:eq-Lambda}.

%% --- Second derivative of p_0 ---

\paragraph{Second derivative of $p_0$.}

\begin{lemma}[Second derivative of $p_0$]\label{sma:lem-d2p0}
  For all $n \ge 1$, $\mu \in (0,1)$, and $\kappa > 0$,
  \begin{equation}\label{sma:eq-d2p0}
    \frac{\partial^2 p_0}{\partial\mu^2}
    = p_0\bigl(\Lambda^2 - \Lambda_2\bigr).
  \end{equation}
\end{lemma}

\begin{proof}
  Starting from $\partial p_0/\partial\mu = -p_0\,\Lambda$, differentiate
  by the product rule:
  \[
    \frac{\partial^2 p_0}{\partial\mu^2}
    = -\frac{\partial p_0}{\partial\mu}\cdot\Lambda
      - p_0\cdot\frac{\partial\Lambda}{\partial\mu}
    = p_0\,\Lambda^2 - p_0\,\Lambda_2
    = p_0\bigl(\Lambda^2 - \Lambda_2\bigr).
  \]
  Here we used $\partial\Lambda/\partial\mu = \Lambda_2$ (since
  $\partial b/\partial\mu = -\kappa$).
\end{proof}

\begin{corollary}[Convexity of $p_0$ in $\mu$]\label{sma:cor-p0-convex}
  For $n \ge 2$:
  $\partial^2 p_0/\partial\mu^2 > 0$, so $p_0$ is strictly convex
  in~$\mu$. For $n = 1$: $p_0 = 1 - \mu$ is linear in~$\mu$.
\end{corollary}

\begin{proof}
  Immediate from~\cref{sma:lem-d2p0} and
  \cref{sma:lem-reciprocals}: for $n \ge 2$,
  $\Lambda^2 > \Lambda_2$ and $p_0 > 0$, so the second derivative is
  strictly positive. For $n = 1$, $\Lambda^2 = \Lambda_2$, giving
  $\partial^2 p_0/\partial\mu^2 = 0$.
\end{proof}

%% --- Second derivative of h ---

\paragraph{Second derivative of $h$.}

\begin{proposition}[Second derivative of $h$]\label{sma:prop-d2h}
  Write $D = 1 - p_0$. Then
  \begin{equation}\label{sma:eq-d2h}
    \frac{\partial^2 h}{\partial\mu^2}
    = \frac{p_0}{D^3}
      \Bigl[-2\Lambda D
            + \mu\Lambda^2(1+p_0)
            - \mu\Lambda_2 D\Bigr].
  \end{equation}
\end{proposition}

\begin{proof}
  From $\partial h/\partial\mu = N/D^2$ where
  $N = \Phi(\mu) = D - \mu\,p_0\,\Lambda$
  (see~\eqref{sma:eq-Phi}), apply the quotient rule:
  \[
    \frac{\partial^2 h}{\partial\mu^2}
    = \frac{N'D - 2N D'}{D^3},
  \]
  where $D' = -\partial p_0/\partial\mu = p_0\,\Lambda > 0$.

  \medskip\noindent\textit{Computing $N'$.}\enspace
  Since $N = (1-p_0) - \mu\,p_0\,\Lambda$, we have
  \[
    N' = p_0\,\Lambda
         - \bigl[p_0\,\Lambda
           + \mu\cdot\tfrac{d}{d\mu}(p_0\,\Lambda)\bigr]
       = -\mu\cdot\frac{d}{d\mu}(p_0\,\Lambda)
       = \mu\,p_0\bigl(\Lambda^2 - \Lambda_2\bigr),
  \]
  using $d(p_0\Lambda)/d\mu = p_0(\Lambda_2 - \Lambda^2)$
  from~\eqref{sma:eq-Phi-deriv}.

  \medskip\noindent\textit{Assembling.}\enspace
  \begin{align*}
    N'D - 2N D'
    &= \mu\,p_0(\Lambda^2 - \Lambda_2)D
       - 2\bigl(D - \mu\,p_0\,\Lambda\bigr)(p_0\,\Lambda) \\
    &= \mu\,p_0(\Lambda^2 - \Lambda_2)D
       - 2D\,p_0\,\Lambda + 2\mu\,p_0^2\Lambda^2 \\
    &= p_0\bigl[-2\Lambda D
       + \mu(\Lambda^2 - \Lambda_2)D
       + 2\mu\,p_0\,\Lambda^2\bigr].
  \end{align*}
  Noting that
  $(\Lambda^2 - \Lambda_2)D + 2p_0\Lambda^2
  = \Lambda^2(D + 2p_0) - \Lambda_2 D
  = \Lambda^2(1 + p_0) - \Lambda_2 D$
  yields the stated form.
\end{proof}

\paragraph{Verification ($n = 1$).}
When $n = 1$, $p_0 = 1 - \mu$, $D = \mu$,
$\Lambda = 1/(1-\mu)$, $\Lambda_2 = \Lambda^2$, and $h \equiv 1$, so
$\partial^2 h/\partial\mu^2 = 0$. In the formula: $N = D - \mu\,p_0\,\Lambda
= \mu - \mu(1-\mu)/(1-\mu) = 0$ identically, so $N' = 0$ and the
numerator vanishes.\enspace$\square$

%% --- Monotonicity of p_0 in all parameters ---

\paragraph{Monotonicity of $p_0$ in all parameters.}

\begin{proposition}[Monotonicity of $p_0$ in $\mu$]\label{sma:prop-p0-mu}
  $\partial p_0/\partial\mu = -p_0\,\Lambda < 0$.
  Increasing $\mu$ strictly decreases the zero probability.
\end{proposition}

\begin{proof}
  Established in~\cref{sma:proof-monotonicity} (auxiliary computation
  preceding the proof of~\cref{thm:monotonicity}).
\end{proof}

\begin{proposition}[Monotonicity of $p_0$ in $\kappa$]\label{sma:prop-p0-kappa}
  For $n \ge 2$ and $\mu \in (0,1)$:
  \begin{equation}\label{sma:eq-dp0-dkappa}
    \frac{\partial p_0}{\partial\kappa}
    = -p_0\,\mu\sum_{j=1}^{n-1}\frac{j}{(b+j)(\kappa+j)} < 0.
  \end{equation}
  For $n = 1$: $\partial p_0/\partial\kappa = 0$ (since
  $p_0 = 1 - \mu$ is independent of~$\kappa$).
\end{proposition}

\begin{proof}
  Differentiating
  $\log p_0 = \sum_{j=0}^{n-1}[\log(b+j) - \log(\kappa+j)]$
  with $\partial b/\partial\kappa = 1-\mu$:
  \[
    \frac{\partial\log p_0}{\partial\kappa}
    = \sum_{j=0}^{n-1}
      \left[\frac{1-\mu}{b+j} - \frac{1}{\kappa+j}\right].
  \]
  For each term, the numerator is
  $(1-\mu)(\kappa+j) - (b+j) = (1-\mu)(\kappa+j) - (1-\mu)\kappa - j = -\mu j$.
  Hence
  \[
    \frac{\partial\log p_0}{\partial\kappa}
    = -\mu\sum_{j=1}^{n-1}\frac{j}{(b+j)(\kappa+j)} < 0,
  \]
  where the $j = 0$ term vanishes. Multiplying by $p_0 > 0$ gives the
  stated formula.
\end{proof}

Increasing $\kappa$ (reducing overdispersion) concentrates mass near
$n\mu > 0$, pulling probability away from zero.

\begin{proposition}[Monotonicity of $p_0$ in $n$]\label{sma:prop-p0-n}
  For all $n \ge 1$, $\mu \in (0,1)$, and $\kappa > 0$:
  \begin{equation}\label{sma:eq-p0-ratio-n}
    \frac{p_0(n+1)}{p_0(n)}
    = \frac{b+n}{\kappa+n} < 1.
  \end{equation}
  Hence $p_0$ is strictly decreasing in~$n$.
\end{proposition}

\begin{proof}
  From the product form~\eqref{eq:p0},
  $p_0(n+1) = p_0(n)\cdot(b+n)/(\kappa+n)$.
  Since $b = (1-\mu)\kappa < \kappa$ for $\mu \in (0,1)$, the ratio
  $(b+n)/(\kappa+n) < 1$.
\end{proof}

\begin{remark}[Polynomial decay of $p_0$ in~$n$]\label{sma:rem-p0-decay}
  The consecutive ratio $(b+n)/(\kappa+n) \to 1$ as $n \to \infty$,
  so $p_0$ decays \emph{polynomially} in~$n$, not exponentially.
  Specifically, by Stirling's approximation for the ratio of rising
  factorials, $p_0(n) \sim C(\mu,\kappa)\,n^{-\mu\kappa}$ as
  $n \to \infty$.  By contrast, the binomial zero probability
  $(1-\mu)^n$ decays exponentially.  This slower decay means that
  even for moderately large group sizes ($n = 50$--$100$), the
  structural-zero probability $p_0$ remains non-negligible when $\kappa$
  is small, reinforcing the importance of the hurdle component
  \citep[see also][for a spatial BB application where this phenomenon
  is empirically relevant]{BandyopadhyayEtAl2011}.
\end{remark}

%% --- Log-derivative and strict monotonicity ---

\paragraph{Log-derivative and strict monotonicity of $\log h$.}

\begin{proposition}[Log-derivative of $h$]\label{sma:prop-logderiv}
  \begin{equation}\label{sma:eq-logderiv}
    \frac{\partial(\log h)}{\partial\mu}
    = \frac{1}{\mu} - \frac{p_0\,\Lambda}{1-p_0}.
  \end{equation}
\end{proposition}

\begin{proof}
  Since $\log h = \log\mu - \log(1-p_0)$, we differentiate each term:
  $d(\log\mu)/d\mu = 1/\mu$, and
  $d\log(1-p_0)/d\mu = (-dp_0/d\mu)/(1-p_0) = p_0\Lambda/(1-p_0)$.
  Subtracting gives the result.
\end{proof}

\begin{theorem}[$\log h$ is strictly increasing]
  \label{sma:thm-logh-increasing}
  For $n \ge 2$:
  \begin{equation}\label{sma:eq-logh-increasing}
    \frac{\partial(\log h)}{\partial\mu} > 0
    \qquad\textup{for all } \mu \in (0,1).
  \end{equation}
\end{theorem}

\begin{proof}
  The condition $\partial(\log h)/\partial\mu > 0$ is equivalent to
  $1/\mu > p_0\Lambda/(1-p_0)$, which rearranges to
  $(1-p_0) > \mu\,p_0\,\Lambda$, i.e.,
  $\Phi(\mu) > 0$ (see~\eqref{sma:eq-Phi}). This was established in
  \cref{thm:monotonicity}.
\end{proof}

\begin{remark}[Decomposition of the log-derivative]
  \label{sma:rem-logderiv-decomp}
  The log-derivative decomposes as
  \[
    \frac{\partial(\log h)}{\partial\mu}
    = \underbrace{\frac{1}{\mu}}_{\text{direct effect}}
    \;-\;
    \underbrace{\frac{p_0\,\Lambda}{1-p_0}}_{\text{truncation attenuation}}.
  \]
  The direct effect represents the proportional increase in the numerator
  $\mu$. The truncation attenuation represents the proportional increase
  in the denominator $1-p_0$. Both terms are positive, but the direct
  effect always dominates because $\Phi(\mu) > 0$
  (\cref{thm:monotonicity}), which in turn follows from the strict
  convexity of~$p_0$ (\cref{sma:cor-p0-convex}).
\end{remark}

%% --- Elasticity ---

\paragraph{Elasticity.}

\begin{definition}[Elasticity of $h$ with respect to $\mu$]
  \label{sma:def-elasticity}
  \begin{equation}\label{sma:eq-elasticity-def}
    \varepsilon_h(\mu)
    = \frac{\partial h}{\partial\mu}\cdot\frac{\mu}{h}
    = \mu\cdot\frac{\partial(\log h)}{\partial\mu}.
  \end{equation}
\end{definition}

\begin{proposition}[Elasticity formula and bounds]
  \label{sma:prop-elasticity-bounds}
  \begin{equation}\label{sma:eq-elasticity-formula}
    \boxed{
      \varepsilon_h(\mu)
      = 1 - \omega(\mu),
      \qquad
      \omega(\mu) = \frac{\mu\, p_0\, \Lambda}{1-p_0}.
    }
  \end{equation}
  Moreover:
  \begin{enumerate}[label=\textup{(\roman*)}]
    \item For $n = 1$: $\varepsilon_h = 0$ (since $h \equiv 1$).
    \item For $n \ge 2$:
      $\varepsilon_h \in (0,1)$ for all $\mu \in (0,1)$.
      The intensity function $h$ is inelastic with respect to~$\mu$.
  \end{enumerate}
  This is the full proof of\/ \cref{prop:elasticity} in the main text.
\end{proposition}

\begin{proof}
  From~\cref{sma:prop-logderiv},
  \[
    \varepsilon_h
    = \mu\left[\frac{1}{\mu} - \frac{p_0\,\Lambda}{1-p_0}\right]
    = 1 - \frac{\mu\, p_0\,\Lambda}{1-p_0}
    = 1 - \omega.
  \]
  By~\cref{sma:thm-logh-increasing}, $\Phi(\mu) > 0$ for $n \ge 2$,
  which gives $\omega < 1$ and hence $\varepsilon_h > 0$.
  Since $\omega > 0$ (all factors are positive), $\varepsilon_h < 1$.
\end{proof}

\begin{remark}[Inelastic interpretation]\label{sma:rem-inelastic}
  A one-percent increase in $\mu$ leads to less than a one-percent
  increase in~$h$, because increasing $\mu$ simultaneously increases
  $1 - p_0$ (the denominator of~$h$), partially offsetting the numerator
  effect.
\end{remark}

\begin{proposition}[Limiting behavior of $\varepsilon_h$]
  \label{sma:prop-elasticity-limits}
  \mbox{}
  \begin{enumerate}[label=\textup{(\alph*)}]
    \item As $\mu \to 0^+$ (for $n \ge 2$):
      $\varepsilon_h \to 0^+$. The truncation attenuation $\omega \to 1^-$.
    \item As $\mu \to 1^-$:
      $\varepsilon_h \to 1 - C_\kappa$, where
      $C_\kappa = \prod_{j=1}^{n-1}j/(\kappa+j)$ is the constant from
      \cref{sma:prop-boundary-one}.  Since
      $p_0 \sim (1-\mu)\,C_\kappa$ and
      $\Lambda \sim 1/(1-\mu) + O(1)$, we get
      $p_0\,\Lambda \to C_\kappa$ and hence
      $\omega \to C_\kappa$.
      For large $n$ or large~$\kappa$, $C_\kappa \approx 0$ and
      the limit approaches~$1$.
    \item As $\kappa \to \infty$ (binomial limit):
      $\varepsilon_h \to 1 - n\mu(1-\mu)^{n-1}\bigl/\bigl[1-(1-\mu)^n\bigr]$.
    \item As $\kappa \to 0$ (maximum overdispersion):
      $p_0 \to 1 - \mu$, $\Lambda \to 1/(1-\mu)$, so $\omega \to 1$ and
      $\varepsilon_h \to 0$. In this limit $h \to 1$; the intensity function
      is flat.
    \item As $n \to \infty$:
      $p_0 \to 0$, so $\omega \to 0$ and $\varepsilon_h \to 1$. The
      truncation correction becomes negligible.
  \end{enumerate}
\end{proposition}

\paragraph{Numerical verification.}
For $n = 10$, $\kappa = 10$, $\mu = 0.3$: the analytic formula gives
$\varepsilon_h = 0.734979$, and centered finite differences yield
$0.734979$. Agreement to six decimal places.\enspace$\square$

%% --- Boundary analysis ---

\paragraph{Boundary analysis.}

\begin{proposition}[Behavior of $h$ as $\mu \to 0^+$]
  \label{sma:prop-boundary-zero}
  For $n \ge 2$ and $\kappa > 0$:
  \begin{equation}\label{sma:eq-boundary-zero}
    h(\mu) \to \frac{1}{\Lambda_0}
    \quad \textup{as } \mu \to 0^+,
  \end{equation}
  where $\Lambda_0 = \kappa\sum_{j=0}^{n-1}1/(\kappa+j)$.
\end{proposition}

\begin{proof}
  Taylor expansion: $p_0(\mu) = 1 - \mu\Lambda_0 + O(\mu^2)$
  (from $p_0(0) = 1$ and $p_0'(0) = -\Lambda_0$). Then
  $1 - p_0 = \mu\Lambda_0 + O(\mu^2)$, and
  $h = \mu\bigl/\bigl(\mu\Lambda_0 + O(\mu^2)\bigr) \to 1/\Lambda_0$.
\end{proof}

For example, when $n = 2$ and $\kappa = 1$:
$\Lambda_0 = 1/1 + 1/2 = 3/2$, so $h(0^+) = 2/3$.
For $n = 10$ and $\kappa = 1$:
$\Lambda_0 = H_{10} \approx 2.929$, so $h(0^+) \approx 0.341$.
As noted in~\cref{rem:boundary}, this limit is strictly less than
unity whenever $n \ge 2$, confirming that the truncation adjustment
has a non-vanishing effect even as $\mu \to 0^+$.

\begin{proposition}[Behavior of $h$ as $\mu \to 1^-$]
  \label{sma:prop-boundary-one}
  For $n \ge 1$ and $\kappa > 0$:
  \begin{equation}\label{sma:eq-boundary-one}
    h(\mu) \to 1 \quad \textup{as } \mu \to 1^-,
    \qquad
    1 - h(\mu) \sim (1-\mu)(1 - C_\kappa),
  \end{equation}
  where $C_\kappa = \prod_{j=1}^{n-1}j/(\kappa+j) \in (0,1)$.
\end{proposition}

\begin{proof}
  As $\mu \to 1$, $b = (1-\mu)\kappa \to 0$, and the product
  form~\eqref{eq:p0} gives $p_0 \sim (1-\mu)\,C_\kappa$ with
  $C_\kappa = \prod_{j=1}^{n-1}j/(\kappa+j)$.
  Then
  \[
    h = \frac{\mu}{1-p_0}
    \sim \frac{1 - (1-\mu)}{1 - (1-\mu)C_\kappa},
  \]
  and
  \[
    1 - h
    \sim \frac{(1-\mu)(1-C_\kappa)}{1-(1-\mu)C_\kappa}
    \sim (1-\mu)(1-C_\kappa)
  \]
  as $\mu \to 1^-$.
\end{proof}

%% --- Conjecture: Monotonicity of elasticity in kappa ---

\paragraph{Elasticity monotonicity in $\kappa$.}

\begin{conjecture}[Monotonicity of $\varepsilon_h$ in $\kappa$]
  \label{sma:conj-elasticity-kappa}
  For fixed $\mu \in (0,1)$ and $n \ge 2$, the elasticity
  $\varepsilon_h$ is increasing in~$\kappa$. Equivalently, the
  attenuation ratio $\omega$ is decreasing in~$\kappa$.
\end{conjecture}

The following proof sketch provides supporting evidence short of a rigorous proof.

\medskip\noindent\textit{Proof sketch.}\enspace
As $\kappa$ increases, $p_0$ decreases (\cref{sma:prop-p0-kappa}),
and the dominant effect on $\omega = \mu\,p_0\,\Lambda/(1-p_0)$ is the
decrease in the numerator factor~$p_0$.
The boundary behavior confirms monotonicity:
$\varepsilon_h \to 0$ as $\kappa \to 0$
(\cref{sma:prop-elasticity-limits}(d)) and
$\varepsilon_h \to 1 - n\mu(1-\mu)^{n-1}/[1-(1-\mu)^n] > 0$ as
$\kappa \to \infty$ (\cref{sma:prop-elasticity-limits}(c)).
A closed-form proof requires showing that the decrease in $p_0$ dominates
the concurrent decrease in~$\Lambda$ uniformly over $(0,1)$.

\begin{remark}[Numerical verification of the conjecture]
  \label{sma:rem-conj-numerical}
  We evaluated $\partial\varepsilon_h/\partial\kappa$ via centered
  finite differences on a grid of $100 \times 100$ values of
  $(\mu,\kappa) \in (0.01, 0.99) \times (0.01, 100)$ for each of
  $n \in \{2, 5, 10, 20, 50, 100\}$. In every case the derivative is
  strictly positive, supporting the conjecture. No counterexample has been
  found.
\end{remark}

%% =============================================================================
%% A.5: Log-Scale Marginal Effect Decomposition
%% Supplements Proposition~\ref{prop:lae-lie} in the main text (Section 2.4)
%% Source: QMD lines 527--608 (03_research-notes-deep-mathematical-properties)
%% Issue 21 resolved: single example replaced with comprehensive threshold table
%% =============================================================================

\subsection{Log-scale marginal effect decomposition}\label{sma:logdecomp}

\cref{prop:lae-lie} established that the marginal effect of covariate~$x_k$
on log expected enrollment decomposes exactly as $\LAE_k + \LIE_k$,
with multipliers $(1-q_i)$ and $(1-\mu_i)\varepsilon_{h,i}$ respectively.
This subsection provides three supplementary results: a correction remark
on the sign structure of the $\LIE$, a sensitivity analysis of the
decomposition ratio with respect to~$\kappa$, and a threshold
characterization of the poverty reversal regime.

%% --- Correction remark on LIE sign ---

\begin{remark}[Sign of the truncation attenuation in the $\LIE$]
  \label{sma:rem-lie-sign}
  The correct $\LIE$ has a \emph{minus} sign in the truncation attenuation:
  \begin{equation}\label{sma:eq-lie-corrected}
    \LIE_k
    = \biggl[\frac{1}{\mu_i}
      - \frac{p_{0,i}\,\Lambda_i}{1-p_{0,i}}\biggr]
      \mu_i(1-\mu_i)\,\tilde{\beta}_{k,s}
    = (1-\mu_i)(1-\omega_i)\,\tilde{\beta}_{k,s},
  \end{equation}
  where $\omega_i = \mu_i\,p_{0,i}\,\Lambda_i/(1-p_{0,i})$ is the
  attenuation ratio from~\cref{sma:prop-elasticity-bounds}. The
  attenuation is \emph{subtracted}, not added, because increasing~$\mu$
  simultaneously increases $1-p_0$ (the denominator of~$h$), which
  attenuates---not amplifies---the response of~$h$ to~$\mu$.
\end{remark}

%% --- Decomposition ratio and sensitivity to kappa ---

\begin{proposition}[Decomposition ratio]\label{sma:prop-DR}
  Define the decomposition ratio
  \begin{equation}\label{sma:eq-DR}
    \textup{DR}(\kappa)
    = \frac{|\LAE_k|}{|\LAE_k| + |\LIE_k|}
    = \frac{(1-q)\,|\tilde{\alpha}_k|}
           {(1-q)\,|\tilde{\alpha}_k|
            + (1-\mu)\,\varepsilon_h\,|\tilde{\beta}_k|}.
  \end{equation}
  Then:
  \begin{enumerate}[label=\textup{(\alph*)}]
    \item As $\kappa \to 0$:
      $\varepsilon_h \to 0$, $\LIE \to 0$, $\textup{DR} \to 1$.
      Maximum overdispersion renders the intensity channel inoperative.
    \item As $\kappa \to \infty$:
      $\varepsilon_h \to \varepsilon_h^{\textup{binom}} \in (0,1)$,
      $\textup{DR} \to \textup{DR}^{\textup{binom}} \in (0,1)$.
    \item Conditional on~\cref{sma:conj-elasticity-kappa},
      $\textup{DR}$ is decreasing in~$\kappa$
      (since $\varepsilon_h$ is increasing in~$\kappa$).
  \end{enumerate}
\end{proposition}

\begin{proof}
  Properties (a) and (b) follow directly from
  \cref{sma:prop-elasticity-limits}(d) and (c), respectively.
  Property~(c) follows because the numerator of~\eqref{sma:eq-DR} is
  constant in~$\kappa$ while the denominator is increasing in~$\kappa$
  through the factor~$\varepsilon_h$.
\end{proof}

\cref{sma:tab-DR} illustrates the sensitivity of the decomposition ratio to
$\kappa$, evaluated at representative parameter values
($n = 10$, $\mu = 0.30$, $q = 0.70$,
$|\tilde{\alpha}| = 0.5$, $|\tilde{\beta}| = 0.8$).

\begin{table}[ht]
\centering
\caption{Sensitivity of the decomposition ratio to~$\kappa$.
  Parameters: $n = 10$, $\mu = 0.30$, $q = 0.70$,
  $|\tilde{\alpha}| = 0.5$, $|\tilde{\beta}| = 0.8$.}
\label{sma:tab-DR}
\smallskip
\begin{tabular}{rcccc}
\toprule
$\kappa$ & $\varepsilon_h$ & $\omega$ & $|\LIE|$ & DR \\
\midrule
0.01  & 0.008 & 0.992 & 0.005 & 0.970 \\
0.1   & 0.071 & 0.929 & 0.040 & 0.790 \\
1     & 0.343 & 0.657 & 0.192 & 0.439 \\
10    & 0.735 & 0.265 & 0.412 & 0.267 \\
100   & 0.859 & 0.141 & 0.481 & 0.238 \\
1000  & 0.874 & 0.126 & 0.489 & 0.235 \\
\bottomrule
\end{tabular}
\end{table}

\begin{remark}\label{sma:rem-DR-sensitive}
  The DR is highly sensitive to $\kappa$ in the range $(0.1, 10)$.
  In our data context, where $\kappa \approx 5$--$15$, both the access and
  intensity channels contribute meaningfully to the total effect.
\end{remark}

%% --- Poverty reversal threshold ---

\begin{proposition}[Threshold for access-effect dominance]
  \label{sma:prop-threshold}
  In the poverty reversal regime
  ($\tilde{\alpha}_{\textup{pov},s} < 0$,
   $\tilde{\beta}_{\textup{pov},s} > 0$),
  the total marginal effect on log expected enrollment is negative
  (access dominates) if and only if
  \begin{equation}\label{sma:eq-threshold}
    \boxed{
      \frac{|\tilde{\alpha}_{\textup{pov},s}|}
           {|\tilde{\beta}_{\textup{pov},s}|}
      > \frac{(1-\mu)\,\varepsilon_h}{1-q}.
    }
  \end{equation}
  The right-hand side is increasing in~$q$ and increasing in~$\kappa$
  (through~$\varepsilon_h$). When participation is high ($q$ near~$1$),
  even a small intensive-margin coefficient can dominate; when $q$ is low,
  the extensive margin dominates more easily.
\end{proposition}

\begin{proof}
  From~\cref{prop:lae-lie}, the total effect is negative iff
  $(1-q)|\tilde{\alpha}| > (1-\mu)\varepsilon_h|\tilde{\beta}|$;
  dividing both sides by $|\tilde{\beta}|(1-q)$ gives
  \eqref{sma:eq-threshold}.
\end{proof}

\cref{sma:tab-threshold} reports the threshold
$(1-\mu)\varepsilon_h/(1-q)$ for parameter combinations spanning the
range of NSECE providers. The elasticity~$\varepsilon_h$ is computed
from the exact formula in~\cref{sma:prop-elasticity-bounds}.

\begin{table}[ht]
\centering
\caption{Access-dominance threshold $(1-\mu)\varepsilon_h/(1-q)$ for
  representative parameter combinations. Values below the empirical
  coefficient ratio $|\alpha_{\textup{pov}}|/|\beta_{\textup{pov}}| = 3.6$
  indicate access dominance. All entries computed with $n = 50$.}
\label{sma:tab-threshold}
\smallskip
\begin{tabular}{llccc}
\toprule
 & & \multicolumn{3}{c}{$\kappa$} \\
\cmidrule(lr){3-5}
$q$ & $\mu$ & 3 & 7 & 15 \\
\midrule
\multirow{3}{*}{0.50}
 & 0.15 & 0.90 & 1.26 & 1.51 \\
 & 0.30 & 1.15 & 1.34 & 1.39 \\
 & 0.45 & 1.04 & 1.10 & 1.10 \\[4pt]
\multirow{3}{*}{0.64}
 & 0.15 & 1.25 & 1.75 & 2.09 \\
 & 0.30 & 1.60 & 1.87 & 1.93 \\
 & 0.45 & 1.44 & 1.52 & 1.53 \\[4pt]
\multirow{3}{*}{0.80}
 & 0.15 & 2.25 & 3.16 & 3.77 \\
 & 0.30 & 2.87 & 3.36 & 3.48 \\
 & 0.45 & 2.60 & 2.74 & 2.75 \\
\bottomrule
\end{tabular}

\medskip
{\footnotesize\textit{Note.} Threshold values are computed from the exact
  elasticity formula~\eqref{sma:eq-elasticity-formula} with
  $p_0 = \prod_{j=0}^{n-1}[(1-\mu)\kappa + j]/(\kappa + j)$
  and $\Lambda = \kappa\sum_{j=0}^{n-1}[(1-\mu)\kappa + j]^{-1}$.
  Values are rounded to two decimal places.}
\end{table}

\begin{remark}\label{sma:rem-threshold-empirical}
  At our posterior estimates
  ($\kappa \approx 7$, $\mu \approx 0.30$, $q \approx 0.64$,
  $n \approx 50$), the threshold is approximately $1.87$
  (\cref{sma:tab-threshold}, row $q = 0.64$, $\mu = 0.30$,
  $\kappa = 7$). Since the empirical coefficient ratio
  $|\alpha_{\textup{pov}}|/|\beta_{\textup{pov}}|
  = 0.324/0.090 \approx 3.6 \gg 1.87$, the access effect dominates
  decisively. For providers near the sample average, the coefficient
  ratio would need to fall below~$1.87$---that is, the extensive-margin
  coefficient would need to be less than twice the intensive-margin
  coefficient---before intensity could dominate. At $q = 0.80$
  (the most extreme participation rate considered), the threshold
  rises to $2.87$--$3.77$, approaching the empirical ratio. This
  confirms that access dominance is most fragile among high-participation
  centers with low~$\mu$ and large~$\kappa$, where the intensity
  elasticity is near unity and the extensive multiplier $(1-q)$ is small.
\end{remark}
